% Partie : sets-functions-relations
% Chapitre : sets
% Section : subsets

\documentclass[../../../include/open-logic-section]{subfiles}

\begin{document}

\olfileid[fr]{sfr}{set}{sub}
\olsection{Sous-ensembles et ensembles des parties}

\begin{explain}
Nous voudrons souvent comparer des ensembles. Une comparaison
naturelle consiste à vérifier si \emph{tout ce qui appartient à
l'un appartient aussi à l'autre}. Cette situation est assez
importante pour justifier une nouvelle notation.
\end{explain}

\begin{defn}[Sous-ensemble]
Si tout !!{element} d'un ensemble $A$ est aussi !!a{element} de~$B$,
on dit que $A$ est un \emph{sous-ensemble}, ou une \emph{partie},
de~$B$, et on écrit $A \subseteq B$. Si $A$ n'est pas un
sous-ensemble de~$B$, on écrit $A \not\subseteq B$.
Si $A \subseteq B$ mais $A \neq B$, on écrit $A \subsetneq B$
et on dit que $A$ est un \emph{sous-ensemble strict} de $B$.
\end{defn}

\begin{ex}
Tout ensemble est un sous-ensemble de lui-même, et $\emptyset$ est
un sous-ensemble de tout ensemble. L'ensemble des entiers naturels
pairs est un sous-ensemble de l'ensemble des entiers naturels.
On a aussi $\{ a, b \} \subseteq \{ a, b, c \}$. En revanche,
si $e$ est distinct de $a$, $b$ et $c$, l'ensemble $\{ a, b, e \}$
n'est pas un sous-ensemble de $\{ a, b, c \}$.\footnote{L'hypothèse
sur $e$, implicite dans l'exemple original, est rendue explicite.}
\end{ex}

\begin{ex}
Le nombre $2$ est un !!{element} de l'ensemble des entiers relatifs,
tandis que l'ensemble des entiers pairs en est un sous-ensemble.
Un ensemble peut toutefois être \emph{à la fois} un !!{element}
et un sous-ensemble d'un autre ensemble : par exemple,
$\{0\} \in \{0, \{0\}\}$ et $\{0\} \subseteq \{0, \{0\}\}$.
\end{ex}

L'extensionnalité donne un critère d'égalité des ensembles :
$A = B$ si et seulement si tout !!{element} de~$A$ est aussi
!!a{element} de~$B$, et réciproquement. La définition de
«~sous-ensemble~» donne à $A \subseteq B$ exactement le sens de
la première moitié de ce critère : tout !!{element} de~$A$ est
aussi !!a{element} de~$B$. Elle s'applique également en échangeant
$A$ et $B$ : $B \subseteq A$ si et seulement si tout !!{element}
de~$B$ est aussi !!a{element} de~$A$. C'est précisément ce
qu'exprime «~et réciproquement~» dans le principe d'extensionnalité.
Autrement dit, ce principe entraîne que deux ensembles sont égaux
si et seulement si chacun est un sous-ensemble de l'autre.

\begin{prop}
$A = B$ si et seulement si $A \subseteq B$ et $B \subseteq A$.
\end{prop}

Introduisons à présent quelques notations utiles. Pour définir
l'inclusion de $A$ dans~$B$, nous avons utilisé l'expression
«~tout !!{element} de~$A$ est \dots~», en remplaçant «~$\dots$~»
par «~!!a{element} de $B$~». Cette \emph{forme}
d'expression est si fréquente qu'une notation formelle nous
sera utile.

\begin{defn}\ollabel{forallxina}
$(\forall x \in A)\phi$ abrège $\forall x(x \in A \lif
\phi)$. De même, $(\exists x \in A)\phi$ abrège $\exists x(x
\in A \land \phi)$.
\end{defn}

Avec cette notation, $A \subseteq B$ si et seulement si
$(\forall x \in A)x \in B$.

Considérons maintenant un type particulier d'ensemble :
l'ensemble de tous les sous-ensembles d'un ensemble donné.

\begin{defn}[Ensemble des parties]
L'ensemble de tous les sous-ensembles d'un ensemble~$A$ est appelé
l'\emph{ensemble des parties de}~$A$ et se note $\Pow{A}$.
  \[
    \Pow{A} = \Setabs{B}{B \subseteq A}
  \]
\end{defn}

\begin{ex}
Quels sont tous les sous-ensembles possibles de $\{ a, b, c \}$ ?
Ce sont $\emptyset$, $\{a \}$, $\{b\}$, $\{c\}$, $\{a, b\}$,
$\{a, c\}$, $\{b, c\}$ et $\{a, b, c\}$. L'ensemble de tous
ces sous-ensembles est $\Pow{\{a,b,c\}}$ :
\[
\Pow{\{ a, b, c \}} = \{\emptyset, \{a \}, \{b\}, \{c\}, \{a, b\},
\{b, c\}, \{a, c\}, \{a, b, c\}\}
\]
\end{ex}

\begin{prob}
Énumérer tous les sous-ensembles de $\{a, b, c, d\}$.
\end{prob}

\begin{prob}
Montrer que, si $A$ a $n$ !!{element}s, alors $\Pow{A}$ a
$2^n$ !!{element}s.
\end{prob}

\end{document}
